\documentclass[11pt,a4paper]{article}
% ---- Préambule commun ----
\usepackage[utf8]{inputenc}
\usepackage[T1]{fontenc}
\usepackage{lmodern}
\usepackage[french]{babel}
\usepackage[a4paper,margin=2.2cm,top=2.6cm,bottom=2.4cm]{geometry}
\usepackage{amsmath,amssymb,amsthm}
\usepackage{enumitem}
\usepackage{xcolor}
\usepackage{listings}
\usepackage{fancyhdr}
\usepackage[most]{tcolorbox}
\usepackage{hyperref}

\definecolor{teal}{HTML}{026573}
\definecolor{tealclair}{HTML}{E6F2F3}
\definecolor{grisclair}{HTML}{F5F5F5}
\definecolor{orangefonce}{HTML}{B35900}

\hypersetup{colorlinks=true,linkcolor=teal,urlcolor=teal,citecolor=teal}

\setlength{\parindent}{0pt}
\setlength{\parskip}{3pt}

% ---- Style Python ----
\lstdefinestyle{python}{
  language=Python,
  basicstyle=\ttfamily\small,
  keywordstyle=\color{teal}\bfseries,
  commentstyle=\color{gray}\itshape,
  stringstyle=\color{orangefonce},
  showstringspaces=false,
  numbers=left,
  numberstyle=\tiny\color{gray},
  numbersep=8pt,
  frame=leftline,
  framerule=1.2pt,
  rulecolor=\color{teal},
  xleftmargin=14pt,
  backgroundcolor=\color{grisclair},
  breaklines=true,
  literate={é}{{\'e}}1 {è}{{\`e}}1 {ê}{{\^e}}1 {à}{{\`a}}1 {ô}{{\^o}}1 {û}{{\^u}}1 {î}{{\^i}}1 {ç}{{\c c}}1
}
\lstset{style=python}

% ---- Compteur d'exercices ----
\newcounter{exo}
\newcommand{\exo}[1]{%
  \stepcounter{exo}%
  \vspace{10pt}%
  \noindent\begin{tcolorbox}[colback=tealclair,colframe=teal,boxrule=1pt,
      left=8pt,right=8pt,top=5pt,bottom=5pt,arc=2pt]
    \textbf{\large Exercice \theexo} \;---\; \textbf{#1}
  \end{tcolorbox}%
  \vspace{4pt}%
  \nopagebreak
}

\newcommand{\partiecours}[1]{%
  \vspace{14pt}%
  {\color{teal}\hrule height 2pt}%
  \vspace{4pt}%
  {\Large\bfseries\color{teal} #1}%
  \vspace{2pt}%
  {\color{teal}\hrule height 0.6pt}%
  \vspace{8pt}%
}

\newtcolorbox{rappel}[1][Rappel]{
  colback=white,colframe=gray!60,boxrule=0.8pt,
  fonttitle=\bfseries,title=#1,coltitle=black,colbacktitle=grisclair,
  left=6pt,right=6pt}

\newtcolorbox{indication}{
  colback=white,colframe=orangefonce!60,boxrule=0.8pt,
  left=6pt,right=6pt,top=4pt,bottom=4pt}

\newcommand{\reponse}{\vspace{4pt}\noindent{\color{teal}\textbf{Solution.}}\;}

\setlength{\headheight}{14pt}
\pagestyle{fancy}
\fancyhf{}
\fancyfoot[C]{\thepage}
\renewcommand{\headrulewidth}{0.4pt}
\renewcommand{\footrulewidth}{0pt}

\fancyhead[L]{\small\color{teal} TD --- Algorithmes de tri}
\fancyhead[R]{\small\color{teal} Corrigé}

\begin{document}

\begin{center}
  {\color{teal}\hrule height 2.5pt}
  \vspace{10pt}
  {\Huge\bfseries\color{teal} Corrigé du TD}\\[6pt]
  {\LARGE\bfseries Les algorithmes de tri}\\[10pt]
  {\large Complexité \;$\bullet$\; Correction \;$\bullet$\; Terminaison}\\[8pt]
  {\normalsize CPGE --- filières MP / PSI \quad$|$\quad Programme officiel marocain 2023}\\[6pt]
  \vspace{6pt}
  {\color{teal}\hrule height 1pt}
\end{center}

\vspace{8pt}

% =====================================================================
\partiecours{Partie I --- Les tris quadratiques}
% =====================================================================

\exo{Tri par sélection : comptage exact et invariant}
\reponse
\textbf{1.} Pour chaque $i\in\{0,\dots,n-2\}$, la boucle interne parcourt
$j\in\{i+1,\dots,n-1\}$, soit $n-1-i$ comparaisons. Au total
\[
  \sum_{i=0}^{n-2}(n-1-i) = \sum_{p=1}^{n-1}p = \frac{n(n-1)}{2}.
\]
Ce nombre ne dépend \emph{pas} du contenu de $t$ : les deux boucles sont des
\texttt{for} de bornes fixées.

\textbf{2.} Un échange par tour de la boucle externe, soit exactement $n-1$ échanges
(y compris les échanges « à vide » lorsque $\texttt{imin}=i$). Meilleur cas comme pire
cas : $\Theta(n^{2})$. Le tri par sélection est le seul des tris étudiés dont la
complexité ne dépend pas de l'entrée.

\textbf{3.} \emph{Terminaison.} Les deux boucles sont des \texttt{for} sur des
intervalles finis : elles terminent. (Formellement, le variant de la boucle interne est
$n-j$, celui de la boucle externe $n-1-i$.)

\textbf{4.} \emph{Invariant} $\mathcal I_i$ : « $t[0..i-1]$ est trié par ordre croissant
et contient les $i$ plus petits éléments du tableau initial (multiensemble) ».

\emph{Initialisation.} Pour $i=0$, $t[0..-1]$ est vide : $\mathcal I_0$ est vraie.

\emph{Hérédité.} Supposons $\mathcal I_i$ vraie. La boucle interne calcule
$\texttt{imin}=\arg\min_{i\le j\le n-1} t[j]$ (invariant interne immédiat :
« $t[\texttt{imin}]=\min t[i..j-1]$ »). Comme $t[0..i-1]$ contient les $i$ plus petits
éléments, $t[\texttt{imin}]$ est le $(i+1)$-ième plus petit élément du tableau initial,
et il est $\ge$ tous les éléments de $t[0..i-1]$. Après l'échange, $t[0..i]$ est donc
trié et contient les $i+1$ plus petits éléments : $\mathcal I_{i+1}$ est vraie.

\emph{Conclusion.} À la sortie, $i=n-1$ et $\mathcal I_{n-1}$ dit que $t[0..n-2]$ est
trié et contient les $n-1$ plus petits éléments ; le dernier élément restant est donc le
maximum et $t$ est trié. De plus seules des transpositions ont été effectuées : le
résultat est bien une permutation de l'entrée.

\textbf{5.} Précisément pour la raison ci-dessus : quand les $n-1$ plus petits éléments
sont placés, le dernier l'est automatiquement. Un tour supplémentaire serait correct mais
inutile.

% ---------------------------------------------------------------------
\exo{Tri par insertion : le rôle des inversions}
\reponse
\textbf{1.} Sur $t=[5,2,4,6,1,3]$ :
\begin{center}\small
\begin{tabular}{|c|c|l|}
\hline
$i$ & \texttt{cle} & tableau après l'itération \\\hline
1 & 2 & $[\,2,5\,|\,4,6,1,3]$ \\
2 & 4 & $[\,2,4,5\,|\,6,1,3]$ \\
3 & 6 & $[\,2,4,5,6\,|\,1,3]$ \\
4 & 1 & $[\,1,2,4,5,6\,|\,3]$ \\
5 & 3 & $[\,1,2,3,4,5,6\,]$ \\\hline
\end{tabular}
\end{center}
(la barre sépare la partie déjà triée du reste.)

\textbf{2.} Si $t$ est trié, le test \texttt{t[j] > cle} est faux dès la première
évaluation : une comparaison par valeur de $i$, soit $n-1$ comparaisons et $0$ décalage.
Complexité $\Theta(n)$.

\textbf{3.} Si $t$ est trié en ordre strictement décroissant, pour chaque $i$ la boucle
\texttt{while} parcourt tout $j=i-1,\dots,0$ puis s'arrête sur $j\ge0$ faux :
$i$ comparaisons et $i$ décalages. Total :
\[
  \sum_{i=1}^{n-1} i = \frac{n(n-1)}{2}\ \text{comparaisons et autant de décalages},
  \qquad \Theta(n^{2}).
\]

\textbf{4.} Considérons un décalage \texttt{t[j+1] = t[j]} : il fait passer l'élément
$t[j]$ (qui vérifie $t[j]>\texttt{cle}$) à droite de \texttt{cle}. Le couple
$(t[j],\texttt{cle})$ était une inversion : il cesse de l'être. Aucun autre couple ne
voit son ordre relatif modifié par ce décalage (les éléments de $t[0..j-1]$ et ceux de
$t[i+1..]$ ne bougent pas, et les éléments déplacés le sont tous d'une case vers la
droite, donc dans le même ordre entre eux). Chaque décalage supprime donc \emph{exactement}
une inversion et n'en crée aucune. Comme le tableau final n'a plus d'inversion, le nombre
total de décalages vaut $I(t)$.

\textbf{5.} Le coût total est : $n-1$ tours de boucle externe à $\Theta(1)$ chacun, plus
$I(t)$ décalages, plus au plus une comparaison « d'arrêt » par tour. D'où
\[
  \Theta\big(n + I(t)\big).
\]
Le tri est donc \emph{adaptatif} : sur un tableau presque trié, $I(t)$ est petit et le
coût est quasi linéaire. Si chaque élément est à distance $\le k$ de sa position finale,
alors $I(t)\le nk$ et la complexité est $\Theta(nk)$ : linéaire à $k$ fixé.
Dans le pire des cas $I(t)=\binom n2$ et l'on retrouve $\Theta(n^2)$.

% ---------------------------------------------------------------------
\exo{Tri par insertion : preuves complètes}
\reponse
\textbf{1.} \emph{Variant} : $V=j+1$. C'est un entier, positif ou nul tant que la
condition $j\ge0$ est vraie, et chaque tour exécute \texttt{j -= 1}, donc $V$ décroît
strictement de $1$. La boucle termine (en au plus $i$ tours).

Quant à l'indice : le \texttt{and} de Python est \emph{paresseux}, il n'évalue
\texttt{t[j] > cle} que si \texttt{j >= 0} est vrai. Sans cela, $j=-1$ donnerait en
Python $t[-1]$ --- le \emph{dernier} élément du tableau --- et l'algorithme serait faux
sans lever d'erreur. C'est un piège classique : l'ordre des deux conditions est
essentiel.

\textbf{2.} \emph{Invariant} $\mathcal J$ : « $t[0..j]$ est inchangé et trié ;
$t[j+2..i]$ contient, dans le même ordre, les éléments initialement en $t[j+1..i-1]$,
et ils sont tous $>\texttt{cle}$ ; la valeur \texttt{cle} est sauvegardée ».

\emph{Initialisation.} Avant le premier tour, $j=i-1$ : $t[j+2..i]=t[i+1..i]$ est vide
et $t[0..i-1]$ est trié par l'invariant externe. \checkmark

\emph{Hérédité.} Si le test est vrai, $t[j]>\texttt{cle}$ ; après
\texttt{t[j+1]=t[j]} puis \texttt{j -= 1}, la tranche $t[j'+2..i]$ (avec $j'=j-1$)
contient exactement l'ancienne tranche augmentée de $t[j]$, en tête, tous
$>\texttt{cle}$, et $t[0..j']$ est inchangé et trié. \checkmark

\emph{Sortie.} La boucle s'arrête soit parce que $j=-1$, soit parce que
$t[j]\le\texttt{cle}$. Dans les deux cas tous les éléments de $t[0..j]$ sont
$\le\texttt{cle}$ (le tableau $t[0..i-1]$ était trié) et tous ceux de $t[j+2..i]$ sont
$>\texttt{cle}$ : l'affectation \texttt{t[j+1] = cle} place \texttt{cle} à sa place et
rend $t[0..i]$ trié.

\textbf{3.} \emph{Invariant} $\mathcal I_i$ : « $t[0..i-1]$ est trié et est une
permutation des $i$ premiers éléments initiaux ». Initialisation : $i=1$, un seul
élément. Hérédité : c'est exactement la conclusion du point \textbf{2}. Sortie : $i=n$,
donc $t[0..n-1]$ est trié et est une permutation de $t$. \checkmark

\textbf{4.} Les seules variables auxiliaires sont \texttt{cle}, \texttt{i} et
\texttt{j} : $\Theta(1)$ mémoire supplémentaire, indépendamment de $n$. Le tri par
insertion est donc \textbf{en place}.

% ---------------------------------------------------------------------
\exo{Tri à bulles et arrêt anticipé}
\reponse
\textbf{1.} \emph{Invariant interne} : « après le tour d'indice $j$, $t[j+1]=\max
t[0..j+1]$ ». Initialisation ($j=0$) : le test place le plus grand de $t[0],t[1]$ en
$t[1]$. Hérédité : si $t[j]=\max t[0..j]$, alors après le test/échange à l'indice $j$,
$t[j+1]=\max(t[j],t[j+1])=\max t[0..j+1]$. \checkmark

\emph{Invariant externe} $\mathcal I_i$ : « après l'itération $i$, la tranche
$t[n-1-i..n-1]$ est triée et contient les $i+1$ plus grands éléments ». En effet la
boucle interne, parcourue jusqu'à $j=n-2-i$, amène le maximum de $t[0..n-1-i]$ en
position $n-1-i$, lequel est $\le$ tous les éléments déjà placés à droite.

\textbf{2.} Variant interne : $(n-1-i)-j$ ; variant externe : $(n-1)-i$.

\textbf{3.} Pire des cas (aucun arrêt anticipé) :
$\sum_{i=0}^{n-2}(n-1-i)=\dfrac{n(n-1)}{2}$ comparaisons, soit $\Theta(n^{2})$.

\textbf{4.} Si $t$ est déjà trié, aucun échange n'a lieu au premier passage :
\texttt{echange} reste \texttt{False} et l'algorithme renvoie après $n-1$ comparaisons.
Complexité $\Theta(n)$ ; la famille est celle des tableaux déjà triés.

\textbf{5.} Un échange n'a lieu que si $t[j]>t[j+1]$, c'est-à-dire sur une inversion de
deux éléments \emph{adjacents} ; l'échange la supprime et ne modifie l'ordre relatif
d'aucun autre couple. Donc le nombre d'échanges vaut exactement $I(t)$ --- même comptage
que pour les décalages du tri par insertion.

\emph{Comparaison.} Les deux tris font le même nombre de « mouvements » $I(t)$, mais pas
le même nombre de \emph{comparaisons}. Le tri à bulles refait un passage complet
($\Theta(n)$ comparaisons) tant qu'un seul échange subsiste. Exemple :
$t=[2,3,4,\dots,n,1]$ a $I(t)=n-1$ inversions ; le tri par insertion coûte $\Theta(n)$,
alors que le tri à bulles a besoin de $n-1$ passages (le $1$ ne remonte que d'une case
par passage) soit $\Theta(n^{2})$. \textbf{Le tri par insertion est nettement préférable
sur un tableau presque trié.}

% ---------------------------------------------------------------------
\exo{Tri en place, tri stable}
\reponse
\textbf{1.} Si les éléments sont de simples nombres, deux éléments égaux sont
indiscernables : échanger leurs positions ne se voit pas. La stabilité n'a de sens que
pour des \emph{enregistrements} (nom, note, classe\dots) triés selon une \emph{clé}
partielle.
\emph{Application.} Pour obtenir une liste d'élèves triée par classe, puis par note
\emph{à l'intérieur de chaque classe}, il suffit de trier d'abord par note, puis par
classe \textbf{avec un tri stable} : le second tri préserve l'ordre des notes au sein
d'une même classe. Avec un tri instable, ce procédé échoue.

\textbf{2.} Soit $t=[2_a,\,2_b,\,1]$.
$i=0$ : $2_b<2_a$ ? non ; $1<2_a$ ? oui, donc $\texttt{imin}=2$. Échange de $t[0]$ et
$t[2]$ : $t=[1,\,2_b,\,2_a]$.
$i=1$ : $2_a<2_b$ ? non, $\texttt{imin}=1$, échange à vide.
Résultat : $[1,\,2_b,\,2_a]$ : $2_b$ précède $2_a$ alors que $2_a$ était initialement
avant. \textbf{Le tri par sélection n'est pas stable.}

\textbf{3.} Dans le tri par insertion, la boucle \texttt{while} s'arrête dès que
$t[j]\le\texttt{cle}$ : la clé est donc insérée \emph{juste après} le dernier élément qui
lui est inférieur ou égal, donc après tous ses égaux déjà placés --- lesquels venaient de
positions plus petites. L'ordre relatif des égaux est préservé : le tri est
\textbf{stable}. Avec \texttt{t[j] >= cle}, la boucle continuerait à décaler les éléments
égaux et \texttt{cle} passerait devant eux : le tri deviendrait instable (tout en restant
correct).

\textbf{4.} Oui : l'échange n'a lieu que si $t[j]>t[j+1]$ (inégalité \emph{stricte}),
donc deux éléments égaux ne sont jamais permutés. Le tri à bulles est \textbf{stable}.

\textbf{5.} Tableau (complété avec les parties II et III) :
\begin{center}\small
\begin{tabular}{|l|c|c|}
\hline
\textbf{Tri} & \textbf{En place} & \textbf{Stable} \\\hline
Sélection & oui & \textbf{non} \\\hline
Insertion & oui & oui \\\hline
Bulles & oui & oui \\\hline
Fusion & \textbf{non} ($\Theta(n)$ auxiliaire) & oui \\\hline
Rapide (Lomuto) & oui ($\Theta(\log n)$ de pile) & \textbf{non} \\\hline
\end{tabular}
\end{center}

% =====================================================================
\partiecours{Partie II --- Le tri par fusion}
% =====================================================================

\exo{La fonction \texttt{fusion} : spécification, invariant, variant}
\reponse
\textbf{1.} \emph{Précondition} : \texttt{A} et \texttt{B} sont deux listes triées par
ordre croissant (au sens large).
\emph{Postcondition} : \texttt{fusion(A,B)} renvoie une \emph{nouvelle} liste triée par
ordre croissant dont le multiensemble des éléments est la réunion de ceux de \texttt{A}
et \texttt{B} ; \texttt{A} et \texttt{B} ne sont pas modifiées.

\textbf{2.} \emph{Variant} : $V=(|A|-i)+(|B|-j)$, entier positif tant que la condition
de boucle est vraie, et chaque tour incrémente exactement l'un de $i$ ou $j$, donc $V$
décroît strictement de $1$. La boucle termine (en au plus $|A|+|B|$ tours).

\textbf{3.} \emph{Invariant} $\mathcal K$ :
\begin{quote}
$C$ est trié ; $C$ contient exactement les éléments de $A[0..i-1]$ et $B[0..j-1]$ ;
et $\max C \le \min\big(A[i..],B[j..]\big)$ (avec la convention $\max\varnothing=-\infty$).
\end{quote}

\emph{Initialisation.} $i=j=0$, $C=[\,]$ : les trois clauses sont vraies. \checkmark

\emph{Hérédité.} Supposons $\mathcal K$ vraie et le test de boucle vrai. Posons
$v=\min(A[i],B[j])$ ; comme $A$ et $B$ sont triées, $v$ est le minimum de tous les
éléments restants, et par $\mathcal K$ il est $\ge\max C$. L'instruction exécutée
(\texttt{C.append(A[i])} si $A[i]\le B[j]$, sinon \texttt{C.append(B[j])}) ajoute
précisément $v$ à la fin de $C$ : $C$ reste trié, la deuxième clause est mise à jour par
l'incrémentation correspondante, et la troisième reste vraie puisque $v$ est retiré des
restes. \checkmark

\emph{Sortie.} La boucle s'arrête quand $i=|A|$ ou $j=|B|$. Supposons $i=|A|$ (l'autre
cas est symétrique) : alors \texttt{A[i:]} est vide et $C + \texttt{B[j:]}$ est trié car
$\max C\le\min B[j..]$ et $B[j..]$ est trié. Les éléments sont exactement ceux de $A$ et
$B$. La postcondition est satisfaite. \checkmark

\textbf{4.} Chaque tour de boucle effectue exactement une comparaison et ajoute
exactement un élément à $C$.
\begin{itemize}[leftmargin=*,itemsep=1pt]
  \item \emph{Minoration.} La boucle ne s'arrête que lorsqu'une des deux listes est
        épuisée, ce qui demande au moins $\min(p,q)$ ajouts : $c\ge\min(p,q)$.
        Borne atteinte si tous les éléments de $A$ sont $\le$ tous ceux de $B$
        (par exemple $A=[1,2]$, $B=[3,4,5]$ : $c=2=\min(p,q)$).
  \item \emph{Majoration.} Après $p+q-1$ ajouts, il ne reste qu'un seul élément, donc
        l'une des deux listes est nécessairement épuisée et la boucle s'est arrêtée :
        $c\le p+q-1$. Borne atteinte pour des listes « en alternance », par exemple
        $A=[1,3,5]$, $B=[2,4,6]$ : $c=5=p+q-1$.
\end{itemize}
Dans tous les cas $c=\Theta(p+q)$, et le coût total (avec les concaténations finales)
est $\Theta(p+q)$.

\textbf{5.} La liste $C$ puis la concaténation occupent $\Theta(p+q)$ emplacements :
\texttt{fusion} n'est \textbf{pas} en place.

% ---------------------------------------------------------------------
\exo{\texttt{tri\_fusion} : terminaison et correction}
\reponse
\textbf{1.} \emph{Fonction de mesure} : $\mu(L)=|L|\in\mathbb N$.
Si $|L|\ge2$, alors $m=\lfloor |L|/2\rfloor$ vérifie $1\le m\le |L|-1$ :
en effet $m\ge1$ car $|L|\ge2$, et $m\le |L|/2<|L|$. Donc
$|L[:m]|=m<|L|$ et $|L[m:]|=|L|-m<|L|$ (car $m\ge1$) : la mesure décroît strictement
sur chacun des deux appels. Comme $\mathbb N$ est bien fondé, la récursion termine.

Si l'on avait écrit \texttt{m = 0}, l'appel \texttt{tri\_fusion(L[0:])} porterait sur
$L$ tout entier : mesure non décroissante, récursion infinie
(\texttt{RecursionError}). Même problème avec \texttt{m = len(L)}.

\textbf{2.} \emph{Récurrence forte sur $n=|L|$.}
\emph{Base} : $n\le1$, la liste est triviallement triée et renvoyée telle quelle.
\emph{Hérédité} : soit $n\ge2$ et supposons la propriété vraie pour toute liste de
longueur $<n$. Comme $|L[:m]|<n$ et $|L[m:]|<n$, \texttt{gauche} et \texttt{droite} sont,
par hypothèse de récurrence, des permutations triées de $L[:m]$ et $L[m:]$.
La précondition de \texttt{fusion} est donc satisfaite, et l'exercice~6 garantit que
\texttt{fusion(gauche,droite)} est une liste triée dont les éléments sont exactement ceux
de $L[:m]$ et $L[m:]$, c'est-à-dire ceux de $L$ (car $L=L[:m]+L[m:]$). \checkmark

\textbf{3.} \textbf{Non}, \texttt{tri\_fusion} n'est pas en place : chaque appel crée
deux tranches puis une liste fusionnée.
À un instant donné, la somme des tailles vivantes le long d'une branche vaut
$n+\frac n2+\frac n4+\dots = \Theta(n)$, et la pile d'appels a une profondeur
$\Theta(\log n)$. L'espace auxiliaire est donc $\Theta(n)$ --- c'est le prix à payer pour
la garantie $\Theta(n\log n)$.

\textbf{4.} Sur $L=[5,2,4,6,1,3]$ :
\begin{center}\small
\begin{tabular}{ll}
\emph{découpage} & $[5,2,4,6,1,3] \to [5,2,4]\ /\ [6,1,3]$\\
                 & $[5,2,4]\to[5]\ /\ [2,4]\to[2]\,/\,[4]$ \\
                 & $[6,1,3]\to[6]\ /\ [1,3]\to[1]\,/\,[3]$ \\[4pt]
\emph{fusions}   & $[2]+[4]\to[2,4]$ ; $[5]+[2,4]\to[2,4,5]$\\
                 & $[1]+[3]\to[1,3]$ ; $[6]+[1,3]\to[1,3,6]$\\
                 & $[2,4,5]+[1,3,6]\to[1,2,3,4,5,6]$\\
\end{tabular}
\end{center}

% ---------------------------------------------------------------------
\exo{Complexité du tri fusion}
\reponse
\textbf{1.} Un appel sur $n\ge2$ éléments découpe en deux listes de tailles
$\lfloor n/2\rfloor$ et $\lceil n/2\rceil$, les trie récursivement, puis les fusionne.
D'après l'exercice~6, la fusion coûte au plus $n-1$ comparaisons (et ce pire cas est
atteignable). D'où
\[
  C(n)=C(\lfloor n/2\rfloor)+C(\lceil n/2\rceil)+(n-1),\qquad C(1)=C(0)=0.
\]

\textbf{2.} Pour $n=2^{k}$ la récurrence devient $C(2^{k})=2C(2^{k-1})+2^{k}-1$.
Posons $D(k)=C(2^{k})$ : $D(k)=2D(k-1)+2^{k}-1$, $D(0)=0$.
\emph{Homogène} : $A\,2^{k}$.
\emph{Particulière pour $2^{k}$} : $2$ étant racine de l'équation caractéristique, on
cherche $\alpha k2^{k}$ :
$\alpha k 2^{k}=2\alpha(k-1)2^{k-1}+2^{k}=\alpha(k-1)2^{k}+2^{k}$ donne $\alpha=1$.
\emph{Particulière pour $-1$} : $c=2c-1$ donne $c=1$.
Donc $D(k)=A2^{k}+k2^{k}+1$ et $D(0)=A+1=0$ donne $A=-1$ :
\[
  D(k)=(k-1)2^{k}+1 \qquad\Longrightarrow\qquad
  \boxed{C(n)=n\log_2 n - n + 1}.
\]
Vérifications : $n=2$ ($k=1$) : $C=0\cdot2+1=1$ (une comparaison pour fusionner deux
singletons) \checkmark ; $n=4$ ($k=2$) : $C=1\cdot4+1=5$ \checkmark.

\textbf{3.} $C(n)=n\log_2 n-n+1=\Theta(n\log n)$.
Dans le \emph{meilleur} des cas, chaque fusion coûte au moins $\min(p,q)=\lfloor n/2\rfloor$
comparaisons, d'où $C_{\min}(n)\ge \frac n2\log_2 n$ : la complexité est encore
$\Theta(n\log n)$. \textbf{Le tri fusion n'est pas sensible à l'ordre initial} : c'est
sa force (garantie) et sa faiblesse (aucun gain sur un tableau presque trié).

\textbf{4.} Montrons $C(n)\le n\lceil\log_2 n\rceil$ par récurrence forte, en posant
$k=\lceil\log_2 n\rceil$. Pour $n=1$, $C(1)=0$. Pour $n\ge2$, on a
$\lceil\log_2\lceil n/2\rceil\rceil = k-1$ et $\lceil\log_2\lfloor n/2\rfloor\rceil\le k-1$,
donc par hypothèse de récurrence
\[
  C(n)\le \Big\lfloor\frac n2\Big\rfloor (k-1)+\Big\lceil\frac n2\Big\rceil (k-1)+n-1
       = n(k-1)+n-1 = nk-1 \le n\lceil\log_2 n\rceil . \qquad\checkmark
\]

\textbf{5.}
\begin{center}\small
\begin{tabular}{|c|c|c|c|}
\hline
$n$ & $n\log_2 n$ & $n^{2}/2$ & rapport \\\hline
$10^{3}$ & $\approx 9{,}97\cdot10^{3}$ & $5\cdot10^{5}$ & $\approx 50$ \\\hline
$10^{6}$ & $\approx 2{,}0\cdot10^{7}$ & $5\cdot10^{11}$ & $\approx 2{,}5\cdot10^{4}$ \\\hline
\end{tabular}
\end{center}
À $10^{8}$ opérations par seconde, $n=10^{6}$ demande $0{,}2$ s au tri fusion contre
environ $1$ heure et $20$ minutes à un tri quadratique.

% ---------------------------------------------------------------------
\exo{Stabilité du tri fusion}
\reponse
\textbf{1.} La liste \texttt{gauche} contient les éléments d'indices initiaux $<m$ et
\texttt{droite} ceux d'indices $\ge m$. Dans \texttt{fusion}, en cas d'\emph{égalité}
le test \texttt{A[i] <= B[j]} est vrai et c'est l'élément de \texttt{A} --- donc celui
d'indice initial le plus petit --- qui est ajouté en premier. Par récurrence forte
(chaque moitié est triée de façon stable), l'ordre relatif des éléments de clés égales
est préservé : \textbf{\texttt{tri\_fusion} est stable}. Tout tient dans le
\texttt{<=}.

\textbf{2.} Avec \texttt{A[i] < B[j]}, l'algorithme reste \textbf{correct} (le résultat
est toujours une permutation triée : en cas d'égalité, l'ordre entre les deux éléments
est indifférent pour le tri). Mais il n'est plus \textbf{stable}.
\emph{Contre-exemple} : $L=[2_a,2_b]$. Alors $m=1$, \texttt{gauche}$=[2_a]$,
\texttt{droite}$=[2_b]$. Le test $2_a<2_b$ est faux (clés égales), donc $2_b$ est
ajouté en premier : le résultat est $[2_b,2_a]$.

\textbf{3.} Avec \texttt{m = 1}, l'algorithme reste correct : les hypothèses de
terminaison ($1\le m\le|L|-1$) et de correction sont toujours vérifiées. Mais la
récurrence devient
\[
  C(n)=C(1)+C(n-1)+\Theta(n)=C(n-1)+\Theta(n)\ \Longrightarrow\ C(n)=\Theta(n^{2}).
\]
On retrouve essentiellement le \textbf{tri par insertion} : à chaque étape on insère un
élément isolé dans une liste déjà triée. Le découpage \emph{équilibré} est donc ce qui
fait tout le gain du tri fusion.

% =====================================================================
\partiecours{Partie III --- Le tri rapide}
% =====================================================================

\exo{\texttt{tri\_rapide} : preuves et cas extrêmes}
\reponse
\textbf{1.} \emph{Mesure} : $\mu(L)=|L|$. Les listes \texttt{petits} et \texttt{grands}
sont construites à partir de \texttt{L[1:]} : c'est le retrait du pivot qui garantit
\[
  |\texttt{petits}|+|\texttt{grands}| = n-1 < n,
\]
donc chacune est de taille $\le n-1<n$. La mesure décroît strictement : la récursion
termine. \emph{Sans} le retrait du pivot (par exemple \texttt{[x for x in L if x <= pivot]}),
la liste \texttt{petits} contiendrait toujours le pivot et pourrait valoir $L$ tout
entier : récursion infinie.

\textbf{2.} \emph{Récurrence forte sur $n=|L|$.} Base : $n\le1$, liste triée.
Hérédité : par hypothèse de récurrence, \texttt{tri\_rapide(petits)} et
\texttt{tri\_rapide(grands)} sont des permutations triées de \texttt{petits} et
\texttt{grands}. Par construction, tout élément de \texttt{petits} est $\le$
\texttt{pivot} et tout élément de \texttt{grands} est $>$ \texttt{pivot} ; la
concaténation
\[
  \underbrace{\text{trié},\ \le \text{pivot}}_{\text{tri\_rapide(petits)}}
  \ \Vert\ [\text{pivot}]\ \Vert\
  \underbrace{\text{trié},\ >\text{pivot}}_{\text{tri\_rapide(grands)}}
\]
est donc triée. Enfin \texttt{petits}, \texttt{[pivot]} et \texttt{grands} forment une
partition du multiensemble $L$ : le résultat est bien une permutation de $L$. \checkmark

\textbf{3.} En comptant une comparaison par élément confronté au pivot, un appel sur $n$
éléments coûte $n-1$ comparaisons (l'implémentation par deux compréhensions en fait
$2(n-1)$, soit un facteur constant $2$). Avec un pivot médian :
\[
  C(n)=(n-1)+2\,C\!\Big(\frac{n-1}{2}\Big).
\]
C'est (à un décalage près) la récurrence du tri fusion : par le théorème général
($a=2$, $b=2$, $\alpha=1$, $f(n)=n$, cas 2) on obtient $C(n)=\Theta(n\log n)$.

\textbf{4.} Si $L$ est trié par ordre croissant, $\texttt{pivot}=L[0]=\min L$, donc
aucun élément de \texttt{L[1:]} n'est $<$ pivot au sens strict : \texttt{petits} est vide
(ou ne contient que les doublons du minimum) et \texttt{grands}$=$\texttt{L[1:]}. D'où
\[
  C(n)=(n-1)+C(n-1),\quad C(0)=0 \quad\Longrightarrow\quad
  C(n)=\sum_{k=1}^{n-1}k=\frac{n(n-1)}{2}=\Theta(n^{2}).
\]
\emph{Paradoxe apparent} : l'entrée la mieux structurée est la pire pour l'algorithme.
Il n'y a pas de contradiction : le tri rapide ne tire sa vitesse que de la
\emph{répartition équilibrée} autour du pivot, et prendre systématiquement le premier
élément d'une liste triée produit le pire déséquilibre possible. C'est un cas fréquent en
pratique (données déjà partiellement ordonnées), d'où l'importance du choix du pivot
(exercice~13).

\textbf{5.} Profondeur de pile : $\Theta(\log n)$ dans le meilleur cas,
$\Theta(n)$ dans le pire. Pour $n=10^{6}$, la version dégénérée provoque un
\texttt{RecursionError} en Python (limite par défaut $\approx1000$), voire un
débordement de pile dans un langage compilé : le problème n'est pas seulement de temps,
mais de \emph{mémoire}.

% ---------------------------------------------------------------------
\exo{Complexité en moyenne du tri rapide}
\reponse
\textbf{1.} Si $L$ est une permutation uniforme de $n$ éléments distincts, $L[0]$ est
uniformément distribué parmi les $n$ éléments : son rang $k$ (nombre d'éléments qui lui
sont strictement inférieurs) suit la loi uniforme sur $\{0,\dots,n-1\}$. Dans ce cas
$|\texttt{petits}|=k$ et $|\texttt{grands}|=n-1-k$, et ces deux sous-listes sont
elles-mêmes des permutations uniformes (propriété classique, admise ici). En
conditionnant sur $k$ et par linéarité de l'espérance :
\[
  C(n)=(n-1)+\frac1n\sum_{k=0}^{n-1}\big(C(k)+C(n-1-k)\big)
      =(n-1)+\frac2n\sum_{k=0}^{n-1}C(k),
\]
la dernière égalité venant du changement d'indice $k\mapsto n-1-k$ dans la seconde somme.

\textbf{2.} En multipliant par $n$ :
$\;n\,C(n)=n(n-1)+2\sum_{k=0}^{n-1}C(k)$.
Au rang $n-1$ :
$\;(n-1)C(n-1)=(n-1)(n-2)+2\sum_{k=0}^{n-2}C(k)$.
Par différence :
\[
  n\,C(n)-(n-1)C(n-1) = (n-1)\big[n-(n-2)\big]+2C(n-1) = 2(n-1)+2C(n-1),
\]
soit $\;\boxed{n\,C(n)=(n+1)\,C(n-1)+2(n-1)}$.

\textbf{3.} En divisant par $n(n+1)$ et en posant $D(n)=C(n)/(n+1)$ :
\[
  D(n)=D(n-1)+\frac{2(n-1)}{n(n+1)} .
\]
Décomposition : $\dfrac{n-1}{n(n+1)}=\dfrac{A}{n}+\dfrac{B}{n+1}$ donne
$n-1=A(n+1)+Bn$, d'où $A=-1$ (en $n=0$) et $B=2$ (en $n=-1$). Ainsi
\[
  \frac{2(n-1)}{n(n+1)} = -\frac2n+\frac{4}{n+1}.
\]

\textbf{4.} Comme $D(1)=C(1)/2=0$, en sommant de $k=2$ à $n$ :
\[
  D(n)=\sum_{k=2}^{n}\Big(-\frac2k+\frac4{k+1}\Big)
      = -2(H_n-1) + 4\Big(H_{n+1}-1-\tfrac12\Big)
      = 4H_{n+1}-2H_n-4 .
\]
Avec $H_{n+1}=H_n+\frac1{n+1}$ :
$\;D(n)=2H_n+\dfrac4{n+1}-4$. Enfin
\[
  C(n)=(n+1)D(n) = 2(n+1)H_n - 4(n+1) + 4
  \qquad\Longrightarrow\qquad \boxed{C(n)=2(n+1)H_n-4n}.
\]
\emph{Vérifications.} $n=2$ : $2\cdot3\cdot\frac32-8=9-8=1$, et deux éléments demandent
bien une comparaison. \checkmark
$n=3$ : $2\cdot4\cdot\frac{11}{6}-12=\frac{44}{3}-12=\frac83\approx2{,}67$ ;
directement $C(3)=2+\frac23\big(C(0)+C(1)+C(2)\big)=2+\frac23=\frac83$. \checkmark

\textbf{5.} $H_n=\ln n+\gamma+o(1)$ donc
\[
  C(n)\sim 2n\ln n = 2\ln 2\cdot n\log_2 n \approx 1{,}386\,n\log_2 n .
\]
Le tri fusion en effectue environ $n\log_2 n$ : le tri rapide fait donc en moyenne
environ $\mathbf{39\,\%}$ de comparaisons en plus.

\textbf{6.} Et pourtant il est souvent plus rapide :
\begin{itemize}[leftmargin=*,itemsep=1pt]
  \item dans sa version en place (exercice~12) il n'alloue aucune mémoire et travaille
        sur des zones contiguës : bien meilleure \emph{localité de cache} ;
  \item la boucle interne de la partition est extrêmement simple (une comparaison, un
        échange conditionnel), donc sa constante multiplicative est petite, alors que la
        fusion recopie chaque élément dans un tableau auxiliaire.
\end{itemize}
Une comparaison de complexités ne dit rien des constantes ni du comportement mémoire :
c'est la limite de l'analyse asymptotique.

% ---------------------------------------------------------------------
\exo{Partition en place (schéma de Lomuto)}
\reponse
\textbf{1.} Sur $t=[2,8,7,1,3,5,6,4]$, $g=0$, $d=7$, $p=t[7]=4$, $i=-1$ :
\begin{center}\small
\begin{tabular}{|c|c|c|c|l|}
\hline
$j$ & $t[j]$ & $t[j]\le4$ ? & $i$ & tableau après le tour\\\hline
0 & 2 & oui & 0 & $[2,8,7,1,3,5,6,4]$ \\
1 & 8 & non & 0 & $[2,8,7,1,3,5,6,4]$ \\
2 & 7 & non & 0 & $[2,8,7,1,3,5,6,4]$ \\
3 & 1 & oui & 1 & $[2,1,7,8,3,5,6,4]$ \\
4 & 3 & oui & 2 & $[2,1,3,8,7,5,6,4]$ \\
5 & 5 & non & 2 & $[2,1,3,8,7,5,6,4]$ \\
6 & 6 & non & 2 & $[2,1,3,8,7,5,6,4]$ \\\hline
\multicolumn{4}{|l|}{échange final $t[3]\leftrightarrow t[7]$} & $[2,1,3,\mathbf4,7,5,6,8]$ \\\hline
\end{tabular}
\end{center}
La fonction renvoie $m=3$.

\textbf{2.} \emph{Invariant} $\mathcal P$ : « $t[g..i]$ ne contient que des éléments
$\le p$ ; $t[i+1..j-1]$ ne contient que des éléments $>p$ ; $t[d]=p$ ; et $t[g..d]$ est
une permutation du contenu initial ».

\emph{Initialisation.} $i=g-1$, $j=g$ : les deux tranches sont vides. \checkmark

\emph{Hérédité.} Deux cas.
\begin{itemize}[leftmargin=*,itemsep=1pt]
  \item $t[j]>p$ : rien n'est modifié, $j$ augmente de $1$ et la tranche
        $t[i+1..j-1]$ s'allonge d'un élément $>p$. \checkmark
  \item $t[j]\le p$ : on incrémente $i$ puis on échange $t[i]$ et $t[j]$.
        Avant l'échange, $t[i]$ (nouveau $i$, donc l'ancien $i+1$) appartenait à la zone
        des $>p$ ; après l'échange, $t[i]=$ ancien $t[j]\le p$ rejoint la zone de gauche
        et l'élément $>p$ est envoyé en position $j$, qui devient la fin de la zone du
        milieu après incrémentation de $j$. \checkmark
\end{itemize}

\emph{Sortie.} $j=d$ : $t[g..i]\le p$, $t[i+1..d-1]>p$ et $t[d]=p$. L'échange final
$t[i+1]\leftrightarrow t[d]$ place le pivot en $m=i+1$, envoie en $t[d]$ un élément
$>p$ (ou $t[d]$ lui-même si $i+1=d$). On obtient bien
\[
  t[g..m-1]\le p=t[m]<t[m+1..d].
\]

\textbf{3.} \texttt{partition} est une boucle \texttt{for} sur un intervalle fini : elle
termine. Pour \texttt{tri\_rapide\_place}, la mesure $\mu(g,d)=d-g+1\in\mathbb N$
décroît strictement : les appels portent sur $(g,m-1)$ et $(m+1,d)$, de mesures
$m-g$ et $d-m$, toutes deux $\le d-g<d-g+1$ puisque $g\le m\le d$. La récursion termine.

\textbf{4.} Exactement $d-g$ comparaisons (une par valeur de $j$ dans
\texttt{range(g,d)}).

\textbf{5.} Chaque appel n'utilise que $\Theta(1)$ variables auxiliaires, mais la pile
d'appels a une profondeur $\Theta(\log n)$ dans le meilleur cas et $\Theta(n)$ dans le
pire. L'algorithme est donc « en place » au sens usuel (aucun tableau auxiliaire), mais
son espace total est $\Theta(\log n)$ en moyenne --- et non $\Theta(1)$ au sens le plus
strict. \emph{Remarque pratique :} rappeler récursivement sur la \emph{plus petite}
moitié et traiter l'autre par une boucle garantit une pile $O(\log n)$ même dans le pire
cas.

\textbf{6.} \textbf{Non, ce tri n'est pas stable.} Sur $t=[2_a,2_b,1]$ avec $g=0,d=2$ :
$p=t[2]=1$, aucun $t[j]\le1$ pour $j\in\{0,1\}$, donc $i$ reste $-1$ et l'échange final
$t[0]\leftrightarrow t[2]$ donne $[1,2_b,2_a]$. Les deux clés égales ont été permutées.
Plus généralement, les échanges de la partition font « sauter » des éléments par-dessus
de grandes distances, ce qui détruit l'ordre initial.

% ---------------------------------------------------------------------
\exo{Choix du pivot et entrées dégénérées}
\reponse
\textbf{1.} Si tous les éléments sont égaux à $v$, alors $p=v$ et le test
\texttt{t[j] <= p} est \emph{toujours} vrai : $i$ est incrémenté à chaque tour et vaut
$d-1$ à la sortie, donc $m=i+1=d$. La partition est maximalement déséquilibrée :
$[g..d-1]$ de taille $n-1$ et $[d+1..d]$ vide. D'où
\[
  C(n)=(n-1)+C(n-1)=\frac{n(n-1)}{2}=\Theta(n^{2}).
\]
Même phénomène pour la version par compréhension : $x\le\texttt{pivot}$ est vrai pour
tous les $x$, donc \texttt{petits}$=$\texttt{L[1:]} et \texttt{grands}$=[\,]$.
\emph{Un tableau constant --- cas fréquent en pratique --- est donc le pire cas du tri
rapide}, ce qui est contre-intuitif.

\textbf{2.} Avec un pivot tiré au hasard, le pire cas $\Theta(n^2)$ \emph{existe encore}
comme \emph{déroulement} possible, mais sa probabilité est infime (il faudrait tomber
systématiquement sur un extremum). Ce qui change est fondamental : il n'y a plus
d'\emph{entrée} mauvaise. La complexité devient une variable aléatoire dont l'espérance
est $\Theta(n\log n)$ \textbf{pour toute entrée}, y compris un tableau déjà trié ; un
adversaire qui connaît l'algorithme ne peut plus construire d'entrée piège.

\textbf{3.} Sur un tableau trié, la médiane de $t[g]$, $t[\lfloor(g+d)/2\rfloor]$,
$t[d]$ est $t[\lfloor(g+d)/2\rfloor]$, c'est-à-dire l'élément \emph{médian} de la
tranche : la partition est parfaitement équilibrée et la complexité devient
$\Theta(n\log n)$. En revanche des entrées quadratiques existent toujours : on peut
construire des suites (dites « tueuses de médiane-de-trois ») pour lesquelles ces trois
positions sont systématiquement extrêmes. Seule la randomisation supprime la notion
d'entrée pire cas.

\textbf{4.} Avec une partition ternaire $(<p)\,(=p)\,(>p)$ et récursion sur les seules
zones extrêmes :
\begin{itemize}[leftmargin=*,itemsep=1pt]
  \item tableau constant : une seule partition en $\Theta(n)$, la zone $(=p)$ absorbe
        tout et il n'y a aucun appel récursif : $\boxed{\Theta(n)}$ ;
  \item tableau à $d$ valeurs distinctes : chaque niveau élimine au moins une valeur et
        la profondeur de l'arbre est $O(\log d)$ en moyenne, d'où $\Theta(n\log d)$ --
        excellent quand $d\ll n$ (notes sur $20$, codes postaux, etc.).
\end{itemize}

% =====================================================================
\partiecours{Partie IV --- Borne inférieure, applications, synthèse}
% =====================================================================

\exo{La borne $\Omega(n\log n)$ : arbre de décision}
\reponse
\textbf{1.} Pour $n=3$ et $t=[a,b,c]$, le tri par insertion compare d'abord $a$ et $b$,
puis insère $c$ (une ou deux comparaisons). L'arbre a $3!=6$ feuilles et une hauteur
de $3$ (branche la plus longue : $c$ plus petit que les deux autres).
Noter que $\lceil\log_2 6\rceil = 3$ : pour $n=3$ ce tri est \emph{optimal}.

\textbf{2.} Soient $\sigma\ne\sigma'$ deux permutations des indices. Il existe une entrée
dont le tri exige de renvoyer $\sigma$, et une autre exigeant $\sigma'$. Deux entrées
conduisant à la même feuille produisent la même suite de réponses aux comparaisons, donc
la même permutation renvoyée. Chaque permutation doit donc être étiquette d'au moins une
feuille : l'arbre a au moins $n!$ feuilles.

\textbf{3.} Par récurrence immédiate sur $h$ : un arbre binaire de hauteur $0$ a $1=2^0$
feuille, et un arbre de hauteur $h$ a au plus deux sous-arbres de hauteur $h-1$, donc au
plus $2\cdot2^{h-1}=2^{h}$ feuilles. Avec le point \textbf{2} : $2^{h}\ge n!$, soit
\[
  h\ \ge\ \log_2(n!).
\]

\textbf{4.} De $n!\ge\left(\frac ne\right)^{n}$ on tire
\[
  \log_2(n!)\ \ge\ n\log_2\frac ne = n\log_2 n - n\log_2 e
  = n\log_2 n - 1{,}4427\,n .
\]
Donc $h=\Omega(n\log n)$ : la hauteur de l'arbre étant le nombre de comparaisons dans le
pire des cas, \textbf{tout tri par comparaisons effectue $\Omega(n\log n)$ comparaisons
dans le pire des cas.}

\textbf{5.} Le tri fusion effectue $n\log_2 n - n + 1$ comparaisons au pire, à comparer
à la borne $n\log_2 n - 1{,}44n$ : il est \textbf{optimal à un terme d'ordre $n$ près},
donc asymptotiquement optimal. Le tri rapide ne l'est pas dans le pire des cas
($\Theta(n^2)$), mais l'est à un facteur constant près ($\approx1{,}39$) en moyenne.

\textbf{6.} Parce que le tri par comptage n'effectue \emph{aucune} comparaison entre
éléments : il utilise leur valeur comme \emph{indice} d'un tableau de compteurs. Il
sort donc du modèle de l'arbre de décision, et la borne ne s'y applique pas : il trie
$n$ entiers de $\{0,\dots,k\}$ en $\Theta(n+k)$, ce qui est linéaire lorsque $k=O(n)$.
\emph{Il n'y a pas de contradiction : la borne ne vaut que dans le modèle des
comparaisons.}

% ---------------------------------------------------------------------
\exo{Trier pour résoudre}
\reponse
\textbf{1.} \emph{Doublons.}
\begin{lstlisting}
def doublons(t):
    u = sorted(t)                 # Theta(n log n)
    for i in range(len(u) - 1):   # Theta(n)
        if u[i] == u[i+1]:
            return True
    return False
\end{lstlisting}
\emph{Correction} : deux éléments égaux sont nécessairement \emph{adjacents} dans le
tableau trié. Version naïve à deux boucles : $\binom n2=\Theta(n^2)$ comparaisons.
Version par tri : $\Theta(n\log n)+\Theta(n)=\Theta(n\log n)$. \emph{Le tri est ici un
prétraitement qui fait tomber la complexité d'un ordre de grandeur.}

\textbf{2.} \emph{Paire de somme $s$.} On trie ($\Theta(n\log n)$) puis on applique la
méthode des deux indices $g=0$, $d=n-1$ : si $t[g]+t[d]<s$ on incrémente $g$, si
$>s$ on décrémente $d$, sinon on a trouvé. Cette phase coûte $\Theta(n)$ (chaque tour
diminue $d-g$ de $1$). Total : $\Theta(n\log n)$, contre $\Theta(n^2)$ pour la version
naïve.

\textbf{3.} \emph{Les $k$ plus grands.}
(i) $k$ tours de tri par sélection : $\sum_{i=0}^{k-1}(n-i)=\Theta(kn)$.
(ii) Tri complet puis lecture des $k$ derniers : $\Theta(n\log n)$.
La première est préférable dès que $kn = O(n\log n)$, c'est-à-dire
$\boxed{k=O(\log n)}$. Pour $k$ grand (par exemple $k=n/2$), le tri complet l'emporte
largement.

\textbf{4.} \emph{Élément majoritaire.} On trie, puis on regarde $u[\lfloor n/2\rfloor]$
et l'on compte ses occurrences en $\Theta(n)$.
\emph{Correction} : si $x$ apparaît $c>n/2$ fois dans le tableau trié, ses occurrences
occupent un bloc d'indices consécutifs $[s,\,s+c-1]$. Si ce bloc ne contenait pas
l'indice $\lfloor n/2\rfloor$, il serait entièrement contenu dans
$[0,\lfloor n/2\rfloor-1]$ ou dans $[\lfloor n/2\rfloor+1,n-1]$, deux intervalles de
longueur $\le n/2$ : contradiction avec $c>n/2$. Donc le majoritaire, s'il existe, est
$u[\lfloor n/2\rfloor]$ ; le comptage final vérifie qu'il existe bien.
Complexité : $\Theta(n\log n)$.

\textbf{5.} \emph{Médiane.} Trier puis lire l'élément central : $\Theta(n\log n)$.
(Il existe des algorithmes de sélection en $\Theta(n)$ --- « médiane des médianes » ---
mais ils sortent du cadre de ce TD.)

% ---------------------------------------------------------------------
\exo{Synthèse et questions de cours}
\reponse
\textbf{1.}
\begin{center}\small
\begin{tabular}{|l|c|c|c|c|c|}
\hline
\textbf{Tri} & \textbf{Meilleur} & \textbf{Moyenne} & \textbf{Pire} &
\textbf{En place} & \textbf{Stable} \\\hline
Sélection & $\Theta(n^2)$ & $\Theta(n^2)$ & $\Theta(n^2)$ & oui & non \\\hline
Insertion & $\Theta(n)$ & $\Theta(n^2)$ & $\Theta(n^2)$ & oui & oui \\\hline
Bulles & $\Theta(n)$ & $\Theta(n^2)$ & $\Theta(n^2)$ & oui & oui \\\hline
Fusion & $\Theta(n\log n)$ & $\Theta(n\log n)$ & $\Theta(n\log n)$ & non & oui \\\hline
Rapide (Lomuto) & $\Theta(n\log n)$ & $\Theta(n\log n)$ & $\Theta(n^2)$ &
oui$^{*}$ & non \\\hline
\end{tabular}
\end{center}
{\footnotesize $^{*}$ aucun tableau auxiliaire, mais $\Theta(\log n)$ de pile en moyenne.}

\textbf{2.} \emph{Vrai ou faux.}
\begin{enumerate}[label=(\alph*),leftmargin=*,itemsep=3pt]
  \item \textbf{Faux.} Les deux notions sont indépendantes : le tri par sélection est en
        place et instable, le tri fusion est stable et non en place.
  \item \textbf{Faux.} $\Theta$ est une propriété \emph{asymptotique} : elle ne dit rien
        pour $n$ petit ni sur les constantes. En pratique le tri par insertion bat le tri
        fusion pour $n\lesssim20$, ce qui justifie les tris hybrides (question \textbf{3}).
  \item \textbf{Faux} tel quel : il faut préciser. Le tri rapide est $\Theta(n\log n)$
        \emph{en moyenne} (et dans le meilleur cas) mais $\Theta(n^2)$ \emph{dans le pire
        des cas}. Une complexité sans quantificateur de cas n'a pas de sens ici.
  \item \textbf{Vrai.} La complexité vaut $\Theta(n+I(t))$ ; si $I(t)=O(n)$ alors
        $\Theta(n+I(t))=\Theta(n)$.
  \item \textbf{Faux}, deux fois. D'une part la borne porte sur le \emph{pire} des cas :
        le tri par insertion n'effectue que $n-1$ comparaisons sur un tableau déjà trié.
        D'autre part elle ne concerne que les tris \emph{par comparaisons} : le tri par
        comptage est en $\Theta(n+k)$.
\end{enumerate}

\textbf{3.} \emph{Tri hybride.}
\begin{enumerate}[label=(\alph*),leftmargin=*,itemsep=3pt]
  \item La correction est préservée : le cas de base « $|L|\le S$ » renvoie une
        permutation triée de $L$ (correction du tri par insertion, exercice~3), ce qui est
        exactement la précondition requise par \texttt{fusion}. La récurrence bien fondée
        de l'exercice~7 s'applique sans changement, avec un cas de base plus large.
  \item L'arbre de découpage est arrêté lorsque les blocs ont une taille $\le S$ : il y a
        donc environ $n/S$ feuilles, chacune triée par insertion en $O(S^{2})$, soit
        $\dfrac nS\cdot O(S^{2})=O(nS)$ pour l'ensemble des feuilles. Au-dessus, il reste
        $\log_2(n/S)$ niveaux de fusions coûtant $\Theta(n)$ chacun. Total :
        \[
          \Theta\!\Big(nS + n\log\frac nS\Big).
        \]
  \item Cette expression reste $\Theta(n\log n)$ tant que $nS = O(n\log n)$, c'est-à-dire
        $S=O(\log n)$. Si $S=\Theta(n)$, le terme $nS$ domine et l'on retombe sur
        $\Theta(n^{2})$ : le tri hybride dégénère en tri par insertion.
  \item Parce que la classe $\Theta$ masque les constantes. Sur un bloc de $10$ à $30$
        éléments, le tri par insertion ne fait aucune allocation, travaille sur une zone
        contiguë de mémoire (bonne localité de cache) et a une boucle interne très
        simple, alors que le tri fusion paie à chaque niveau la création de tranches et
        la recopie. C'est exactement la stratégie utilisée par les bibliothèques
        standard (\emph{Timsort} en Python, \emph{introsort} en C++).
\end{enumerate}

\vspace{14pt}
{\color{teal}\hrule height 1pt}
\vspace{6pt}
\begin{center}
\small\itshape Fin du corrigé --- 16 exercices.
\end{center}

\end{document}
